% !PRINTED 486
\weberpage{486}
%<<CONT
\[
\begin{aligned}
y &= \psi(\alpha, \alpha_1 \ldots \alpha_{r-1}) \\
y_1 &= \psi(\beta, \beta_1 \ldots \beta_{r-1}) \\
& \phantom{=} \; \cdots\cdots\cdots\cdots \\
y_{s-1} &= \psi(\sigma, \sigma_1 \ldots \sigma_{r-1})
\end{aligned}
\tag{8}
\]
are all different from one another.

To see the possibility of this, we can, for example, set
\[
\begin{aligned}
\psi(t, A) &= (t - \alpha)(t - \alpha_1) \ldots (t - \alpha_{r-1}), \\
\psi(t, B) &= (t - \beta)(t - \beta_1) \ldots (t - \beta_{r-1}), \\
& \phantom{=} \; \cdots\cdots\cdots\cdots \\
\psi(t, S) &= (t - \sigma)(t - \sigma_1) \ldots (t - \sigma_{r-1})
\end{aligned}
\tag{9}
\]
and then (by \S.\,143, 1.) a rational value can be found for $t$ such
that these quantities all turn out different from one another. These
values can therefore be taken for the $y$ in (8).

If we now understand by $u$ an independent variable and set
\weq{10}{\varphi(u) = (u - y)(u - y_1) \ldots (u - y_{s-1})}
then $\varphi(u)$, since its coefficients remain unchanged under the permutations
of the group, is a function in $\Omega$ whose roots are the quantities $y$. \gesperrt{It is moreover irreducible}, for
from the assumed transitivity of the group it follows that
if a rational function of $y$ vanishes for one of the values (8),
it must also vanish for all the others.

If $\omega$ is any symmetric function of the $\alpha, \alpha_1 \ldots \alpha_{r-1}$,
then from the consideration of the expression
\[
\varphi(u) \left( \frac{\omega}{u - y} + \frac{\omega_1}{u - y_1} + \cdots + \frac{\omega_{s-1}}{u - y_{s-1}} \right),
\]
which is an integral rational function of $u$ in $\Omega$, we conclude, just as above,
by setting $u = y$, that $\omega$ can be represented rationally
in terms of $y$; and accordingly the function
\weq{11}{\psi(t, A) = \psi(t, y),}
whose roots are the quantities $\alpha, \alpha_1 \ldots \alpha_{r-1}$, can also be expressed
rationally in terms of $y$.

\gesperrt{The equation $\psi(t, A) = 0$ is also irreducible
in the field $\Omega(y)$}; for if $\psi_1(t, y)$ is a rational divisor
of $\psi(t, y)$, and if
%CONT>>
